\par\Needspace{6\baselineskip}\medskip
\begin{center}\hypertarget{native-187-article140}{}\pdfbookmark[2]{140.}{article-140}140.\end{center}
Primo observamus, si $N$ fuerit punctum quodcunque circuli maximi $CC'C''$,
distantiaeque punctorum $C,C',C''$ a puncto $N$ secundum directionem eandem
numerentur, quae tendit a $C$ ad $C''$, ita ut generaliter fiat
\[
NC''-NC'=2f,\qquad NC''-NC=2f',\qquad NC'-NC=2f'',
\]
haberi aequationem
\[
\text{I.}\qquad 0=\sin2f\sin NC-\sin2f'\sin NC'+\sin2f''\sin NC''.
\]
Iam supponemus, $N$ accipi in intersectione circulorum maximorum $BB^*B''$,
$CC'C''$, quasi in nodo ascendente prioris supra posteriorem. Designemus per
$\mathfrak{C},\mathfrak{C}',\mathfrak{C}'',\mathfrak{D},\mathfrak{D}',\mathfrak{D}''$ resp. distantias punctorum $C,C',C'',D,D',D''$ a circulo
maximo $BB^*B''$, ab alterutra ipsius parte positive, ab altera opposita negative
acceptas. Hinc manifesto $\sin\mathfrak{C},\sin\mathfrak{C}',\sin\mathfrak{C}''$ resp. proportionales erunt ipsis
$\sin NC,\sin NC',\sin NC''$, unde aequatio I sequentem induit formam
\[
0=\sin2f\sin\mathfrak{C}-\sin2f'\sin\mathfrak{C}'+\sin2f''\sin\mathfrak{C}''
\]
sive multiplicando per $rr'r''$
\hypertarget{native-188}{}
\[
\text{II.}\qquad 0=nr\sin\mathfrak{C}-n'r'\sin\mathfrak{C}'+n''r''\sin\mathfrak{C}''.
\]
Porro patet, esse $\sin\mathfrak{C}$ ad $\sin\mathfrak{D}'$ ut sinum distantiae puncti $C$ a $B$ ad sinum
distantiae puncti $D'$ a $B$, utraque distantia secundum eandem directionem mensurata. Habetur itaque
\[
-\sin\mathfrak{C}=\frac{\sin\mathfrak{D}'\sin CB}{\sin(AD'-\delta)}\mathpunct{\raisebox{.45ex}{,}}
\]
prorsusque simili modo eruitur
\[
\begin{aligned}
-\sin\mathfrak{C}&=\frac{\sin\mathfrak{D}''\sin CB}{\sin(AD''-\delta)}\\[.8ex]
-\sin\mathfrak{C}'&=\frac{\sin\mathfrak{D}\sin C'B^*}{\sin(A'D-\delta'+\sigma)}
=\frac{\sin\mathfrak{D}''\sin C'B^*}{\sin(A'D''-\delta'+\sigma)}\\[.8ex]
-\sin\mathfrak{C}''&=\frac{\sin\mathfrak{D}\sin C''B''}{\sin(A''D-\delta'')}
=\frac{\sin\mathfrak{D}'\sin C''B''}{\sin(A''D'-\delta'')}\mathpunct{\cdot}
\end{aligned}
\]
Dividendo itaque aequationem II per $r''\sin\mathfrak{C}''$, prodit
\[
\begin{aligned}
0={}&n\cdot\frac{r\sin CB}{r''\sin C''B''}\cdot\frac{\sin A''D'-\delta''}{\sin(AD'-\delta)}\\[.8ex]
&-n'\cdot\frac{r'\sin C'B^*}{r''\sin C''B''}\cdot\frac{\sin(A''D-\delta'')}{\sin(A'D-\delta'+\sigma)}+n''.
\end{aligned}
\]
Quodsi hic arcum $C'B'$ per $z$ designamus, pro $r,r',r''$ valores suos ex art.
praec. substituimus, brevitatisque caussa ponimus
\[
\begin{aligned}
11)\quad&\frac{R\sin\delta\sin(A''D'-\delta'')}{R''\sin\delta''\sin(AD'-\delta)}=a\\[.8ex]
12)\quad&\frac{R'\sin\delta'\sin(A''D-\delta'')}{R''\sin\delta''\sin(A'D-\delta'+\sigma)}=b,
\end{aligned}
\]
aequatio nostra ita se habebit
\[
\text{III.}\qquad 0=an-bn'\cdot\frac{\sin(z-\sigma)}{\sin z}+n''.
\]
Coëfficientem $b$ etiam per formulam sequentem computare licet, quae ex aequationibus modo allatis facile deducitur:
\[
13)\qquad a\cdot\frac{R'\sin\delta'\sin(AD''-\delta)}{R\sin\delta\sin(A'D''-\delta'+\sigma)}=b.
\]
Calculi confirmandi caussa haud inutile erit, utraque formula 12 et 13 uti.\HistoricalUnmarkedNote{[Handschriftliche Bemerkung:] Ad fig. 4: Sit $P$ polus circuli maximi $CC'C''$ eritque
\[
\begin{aligned}
&\left(\frac{n''\cos BP}{N''\sin\beta}-1\right)\sin A'D''.\sin B'D\\[.7ex]
&\qquad=\left(\frac{n\cos B''P}{N\sin\beta''}-1\right)\sin A'D.\sin B'D''
\end{aligned}
\]
[wo $N,N''$ die doppelten Dreiecksflächen zwischen den Radien vectoren der Erde sind.]}
\hypertarget{native-189}{}Quoties $\sin(A'D''-\delta'+\sigma)$ maior est quam $\sin(A'D-\delta'+\sigma)$, formula posterior
a tabularum erroribus inevitabilibus minus afficietur, quam prior, adeoque huic
praeferenda erit, si forte parvula discrepantia illinc explicanda in valoribus
ipsius $b$ se prodiderit; contra formulae priori magis fidendum erit, quoties
$\sin(A'D''-\delta'+\sigma)$ minor est quam $\sin(A'D-\delta'+\sigma)$: si magis placet, medium
idoneum inter ambos valores adoptabitur.

Calculo examinando sequentes quoque formulae inservire possunt, quarum
tamen derivationem non ita difficilem brevitatis caussa supprimimus:
\[
\begin{aligned}
0&=\frac{a\sin(l''-l')}{R}-\frac{b\sin(l''-l)}{R'}\cdot\frac{\sin(\delta'-\sigma)}{\sin\delta'}+\frac{\sin(l'-l)}{R''}\\[.8ex]
b&=\frac{R'\sin\delta'}{R''\sin\delta''}\cdot\frac{U\cos\beta\cos\beta''}{\sin(AD'-\delta)\sin\epsilon'}\mathpunct{\raisebox{.45ex}{,}}
\end{aligned}
\]
ubi $U$ exprimit quotientem $\frac{S}{\sin(\delta'-\sigma)}=\frac{T\sin(t+\gamma')}{\cos(\delta'-\sigma)}$ (art. 138 aequ. 10).
